\section{Design}\label{design}

This chapter describes the high-level design adopted to satisfy the requirements in \cref{requirements}. The main design goal is to support autonomous nodes (FR1), local sensor ingestion and replicated convergence (FR2--FR4), decentralized membership and liveness observation (FR5--FR6), recovery after failures (FR7), and read-only monitoring (FR8, FR12), while preserving decentralization, eventual consistency, partition tolerance, and course-scale experimental scalability (NFR1--NFR3, NFR7).

\subsection{Architecture}\label{architecture}

The system follows a \textbf{fully decentralized peer-to-peer architecture} in which every node executes the same logical stack and can both produce local sensor data and consume data generated by other peers. There is no master node, no dedicated coordinator, and no centralized database. This architectural style directly addresses FR1 and NFR1: any node can participate in the cluster as an autonomous peer, and the failure of one peer does not invalidate the whole system.

From a logical viewpoint, the architecture is organized into three cooperating planes:
\begin{itemize}
  \item a \textbf{control plane}, responsible for decentralized peer discovery, membership maintenance, peer knowledge propagation, adaptive liveness observation, Phi Accrual status classification, and topology knowledge;
  \item a \textbf{data plane}, responsible for local sensor ingestion, local state maintenance, deterministic merge semantics, push-pull anti-entropy exchange, and full-state rejoin synchronization;
  \item an \textbf{observability plane}, responsible for exposing cluster snapshots, fault evidence, metrics, and monitoring views to operators and observer clients.
\end{itemize}

This separation is important because the project intentionally distinguishes between \emph{knowing who is in the system} and \emph{replicating what the system knows about sensors}. Membership and liveness information evolve under different timing assumptions and serve different purposes than sensor-state replication. Keeping the two concerns separate improves modularity, clarity, and evolvability, which supports NFR6 and the maintainability goals discussed in \cref{ds-features}.

Within each node, the architecture is further layered:
\begin{itemize}
  \item a \textbf{transport layer}, dealing only with network communication;
  \item a \textbf{protocol layer}, defining message types and dispatching received messages to the correct logic;
  \item a set of \textbf{domain components}, where the actual semantics of membership, failure detection, replication, topology, and observation are applied.
\end{itemize}

This layered organization keeps low-level communication concerns separate from distributed-system semantics. As a consequence, the same conceptual design would still hold if the underlying transport or encoding technology changed, which is consistent with the intended role of the design chapter.
\subsection{Infrastructure}\label{infrastructure}

The infrastructure is intentionally minimal. The essential infrastructural unit is the \textbf{node process}, replicated one time for each participant in the cluster. Every node contains the same internal subsystems: sensor ingestion, local state management, peer membership handling, failure detection, TCP protocol handling, and a read-only HTTP observation interface. This homogeneous-node model simplifies deployment and supports both FR1 and NFR5, because the same runtime can be instantiated in different topologies without role specialization.

No external infrastructural component is required for normal operation. In particular, the design does not rely on:
\begin{itemize}
  \item a central database;
  \item a message broker;
  \item a coordination service;
  \item a load balancer;
  \item a dedicated monitoring backend.
\end{itemize}

The only optional external-facing component is the browser dashboard, which is not part of the distributed core and behaves as a read-only observer. It can be hosted separately because it only consumes the observation API exposed by any node.

The distributed deployment model is therefore the following:
\begin{itemize}
  \item one or more nodes can run as separate processes, containers, or machines, while distributed behavior is exercised when at least two nodes are active;
  \item each node may manage a different local set of sensor sources or sensor providers;
  \item nodes communicate directly with peers over the inter-node network;
  \item observers may connect to any node's read-only HTTP API, or use the dashboard configured to consume that API.
\end{itemize}

Component discovery is bootstrapped through a configured set of initial peers. After first contact, discovery is no longer purely static: nodes can exchange peer knowledge and extend their local membership view through join messages, peer-list exchange, membership gossip, and full-state synchronization metadata. Naming is based on stable node identifiers plus network coordinates required for communication. This choice is sufficient for the project scope because it avoids introducing centralized service-discovery infrastructure while still allowing evolving cluster knowledge, which is coherent with FR5 and NFR1.

From a physical standpoint, the design does not assume geographic distribution at continental scale; however, it does assume that nodes may be distributed across separate machines or containers and may suffer message delay, temporary unreachability, or partial connectivity. This is exactly the failure model needed to study NFR2 and NFR3.

\subsection{Modelling}\label{modelling}

The main domain entities are:
\begin{itemize}
  \item \textbf{Node}, representing an autonomous distributed participant;
  \item \textbf{Peer}, representing another known node together with its membership and liveness metadata;
  \item \textbf{Sensor source}, representing a local producer of readings;
  \item \textbf{Sensor reading}, representing one observation generated by a sensor source;
  \item \textbf{Replicated sensor record}, representing the current winning value known for one logical sensor identifier;
  \item \textbf{Membership entry}, representing a node's local belief about another peer's status;
  \item \textbf{Topology entry}, representing adjacency knowledge useful for network observability and connection policies;
  \item \textbf{Observer view}, representing an aggregated, read-only projection of internal distributed state.
\end{itemize}

These entities map naturally to infrastructural components. Sensor sources are local to a single node, while replicated sensor records, membership entries, and topology entries are stored as node-local replicas. In other words, the design does not centralize domain state in one authoritative repository; instead, each node stores its own best-effort replica of the distributed knowledge relevant to operation and observation.

The most important domain events are:
\begin{itemize}
  \item production of a local sensor reading;
  \item discovery of a new peer;
  \item reception of a heartbeat;
  \item transition of a peer between liveness states;
  \item emission or reception of a sensor-state update;
  \item request for missing replication data;
  \item full-state transfer during recovery;
  \item generation of an introspection snapshot for observers.
\end{itemize}

The exchanged messages can be grouped into three categories:
\begin{itemize}
  \item \textbf{control messages}, for discovery, membership dissemination, and liveness observation;
  \item \textbf{state-replication messages}, for sensor update propagation and catch-up;
  \item \textbf{read-only queries}, for HTTP observation by users or tools.
\end{itemize}

The system state includes:
\begin{itemize}
  \item latest winning sensor values together with origin and timestamp metadata;
  \item membership knowledge about known peers;
  \item liveness evidence and peer-status information;
  \item topology knowledge and connection-related views;
  \item operational events and metrics used for observability.
\end{itemize}

This model is intentionally centered on \emph{current replicated knowledge}, not on long-term historical storage. That design keeps the project focused on convergence, liveness, and fault handling rather than on archival analytics.

\subsection{Interaction}\label{interaction}

There are two main interaction families in the system.

The first is \textbf{peer-to-peer interaction among nodes}. This interaction is best-effort, message-based, and decoupled through protocol dispatch. It is used for:
\begin{itemize}
  \item bootstrap and peer discovery;
  \item heartbeat exchange;
  \item dissemination of membership knowledge;
  \item dissemination and retrieval of replicated sensor state;
  \item recovery synchronization after missed updates.
\end{itemize}

The second is \textbf{observer interaction}, where external clients query a node through the HTTP read API. This interaction is request-response and read-only. It does not participate in cluster coordination and it never mutates the distributed state.

At the inter-node level, communication follows a layered path:
\begin{quote}
\emph{domain intent} $\rightarrow$ \emph{protocol message} $\rightarrow$ \emph{transport delivery} $\rightarrow$ \emph{protocol dispatch} $\rightarrow$ \emph{domain update}
\end{quote}

The main interaction patterns are:
\begin{itemize}
  \item \textbf{bootstrap and peer-list exchange}, used when a node joins and needs initial peer knowledge;
  \item \textbf{periodic heartbeat exchange}, used to collect liveness evidence;
  \item \textbf{gossip dissemination}, used to spread membership knowledge without central coordination;
  \item \textbf{push-pull state synchronization}, used to exchange recent sensor-state updates efficiently;
  \item \textbf{full synchronization on demand}, used when incremental recovery is insufficient;
  \item \textbf{polling-based observation}, used by the dashboard and external tools to inspect cluster state.
\end{itemize}

This design supports FR3, FR5, FR6, FR7, FR8, FR9, FR10, FR11, and FR12 because it combines low-coupling peer communication with an observation interface that stays orthogonal to the replicated core.

\subsection{Behaviour}\label{behaviour}

The behavior of the main components can be summarized as follows.

\textbf{Sensor sources} are active components: they periodically generate local readings and inject them into the node. They are conceptually producers and do not need global knowledge.

\textbf{The local state manager} is stateful and authoritative for the node's replica. It receives local and remote updates, applies deterministic merge rules, stores the current winners, and exposes snapshots to both replication and observability components. This component is central to FR2, FR3, and FR4.

\textbf{The membership manager} is also stateful. It stores known peers, updates direct or indirect liveness evidence, and maintains the node's local view of cluster membership.

\textbf{The failure detector} behaves as a local evaluator rather than as a replicated authority. It interprets heartbeat timing and produces suspicion levels for peers observed directly by the node. Those local judgements are then translated into membership information that can be disseminated.

\textbf{The heartbeat and gossip runtime} is periodic. At each round, it evaluates liveness, sends heartbeat probes, and disseminates updated membership knowledge. Its purpose is not to guarantee exact membership synchronization at every instant, but to make peer-status knowledge progressively converge.

\textbf{The replication runtime} is periodic as well. It pushes recent sensor-state updates to selected peers and pulls missing updates from others. When incremental recovery is not sufficient, it can trigger a broader synchronization phase. This behavior is essential for FR7 and NFR2.

\textbf{The observability subsystem} is read-only. It collects snapshots from state, membership, topology, and event/metric providers, then publishes a coherent observation view for users and tools.

State updates are therefore concentrated in a small set of stateful components:
\begin{itemize}
  \item local sensor-state storage;
  \item local membership storage;
  \item local topology and observability stores.
\end{itemize}

Other components mainly trigger transitions or move data between these stores. This separation keeps write paths explicit and supports maintainability.

\subsection{Data and Consistency Issues}\label{data-and-consistency-issues}

The system stores several forms of node-local runtime data:
\begin{itemize}
  \item the replicated sensor-state winners;
  \item the incremental replication view needed for propagation and recovery;
  \item the membership table and liveness metadata;
  \item topology knowledge;
  \item recent events and metrics for introspection.
\end{itemize}

This data is primarily \textbf{runtime state}, not long-term persistent storage. The design does not require an external database because the goal is to study distributed replication and convergence of current system knowledge rather than durable business records. This keeps the architecture lightweight and self-contained, matching IR1.

Conceptually, the replicated sensor state behaves like a key-value map from a logical sensor identifier to the current winning record. The important property is not relational querying power, but deterministic merge semantics. For this reason, the core consistency problem is solved at the application level rather than delegated to a storage engine.

Consistency is handled differently for different data types:
\begin{itemize}
  \item \textbf{sensor state} uses eventual consistency with deterministic Last-Writer-Wins conflict resolution based on timestamp and origin metadata;
  \item \textbf{membership state} also converges eventually, but it reflects local observations that may temporarily differ across nodes;
  \item \textbf{observability data} is inherently approximate and time-sensitive, because it represents a snapshot of changing distributed state.
\end{itemize}

This means the system explicitly rejects strong consistency as a primary goal. During partitions or transient delays, different nodes may expose different but still valid local views. Once communication resumes and no newer conflicting updates appear, the replicas converge. This choice is aligned with NFR2 and NFR3.

Data is shared among components in well-defined ways:
\begin{itemize}
  \item sensor producers share generated readings with the local state manager;
  \item the state manager shares snapshots and deltas with the replication runtime and the observability subsystem;
  \item the membership subsystem shares peer-status information with heartbeat, gossip, topology, and observability components;
  \item the observability subsystem reads from other stores but does not modify them.
\end{itemize}

This read/write asymmetry is deliberate: observation must remain decoupled from mutation so that FR8 can be satisfied without introducing accidental feedback paths.

\subsection{Fault-Tolerance}\label{fault-tolerance}

Fault tolerance is one of the central design drivers of the project. The first mechanism is \textbf{state replication}: sensor-state knowledge is disseminated among peers so that the system does not depend on a unique owner of cluster knowledge. Replication is best-effort and decentralized, which supports NFR1 and NFR3 even though it does not provide instantaneous consistency.

The second mechanism is \textbf{heartbeat-based liveness observation}. Nodes periodically exchange liveness probes and evaluate the timing of received heartbeats. This allows each node to form a local judgement about whether a peer is alive, suspected, or dead. Since such knowledge is local and timing-dependent, the design treats liveness as a continuously revised estimate rather than as an absolute truth.

The third mechanism is \textbf{membership gossip}. A node does not keep liveness knowledge only for itself; it disseminates updated membership information so that peer knowledge can propagate even when direct connectivity patterns are incomplete. This improves resilience of membership awareness and helps the system maintain a broader view of the cluster.

The fourth mechanism is \textbf{anti-entropy recovery}. Incremental propagation is efficient when peers stay mostly synchronized, but it may be insufficient after outages, restart, or missed traffic. For this reason, the design includes a stronger recovery path that allows a lagging node to regain a coherent replicated view after disruption, directly supporting FR7.

Error handling is intentionally failure-aware rather than failure-oblivious. If a component or peer becomes temporarily unavailable:
\begin{itemize}
  \item surviving nodes continue local operation;
  \item local replicas may diverge temporarily;
  \item membership views reflect degraded reachability;
  \item synchronization resumes when communication becomes possible again.
\end{itemize}

This behavior is preferable to halting the whole system because the project prioritizes availability and observation under faults over strong global coordination.

\subsection{Availability}\label{availability}

The design does not introduce classical web-style caching layers or external load balancers, because the system is not centered on high-volume client traffic toward a centralized service. Instead, availability is achieved structurally through replication, autonomy of nodes, and local read access to each node's current view.

Each node acts as a locally available observation point. As long as a node remains running, it can still:
\begin{itemize}
  \item ingest its own local sensor readings;
  \item maintain its local replica;
  \item expose read-only snapshots to operators and tools.
\end{itemize}

In this sense, replicated state itself plays a role analogous to availability-oriented caching: each node keeps a local working copy of shared cluster knowledge, so read operations do not require contacting a central authority.

No explicit load balancing is required because there is no single mandatory ingress point for observation. Different observers may query different nodes, and each node exposes its own local view. This is sufficient for the experimental scope of the project.

During a \textbf{network partition}, the design favors continued local operation over blocking. Nodes inside each connected component continue to exchange information among themselves, while nodes separated by the partition stop receiving fresh updates from one another. Consequently:
\begin{itemize}
  \item sensor-state replicas may diverge temporarily across partitions;
  \item membership views may mark remote peers as suspected or dead;
  \item local observation remains available inside each connected component.
\end{itemize}

When the partition heals, incremental dissemination and recovery mechanisms progressively re-align the replicas. This behavior is a direct consequence of choosing availability and partition tolerance over strong consistency, and it is exactly the semantics required by NFR2 and NFR3.

\subsection{Security}\label{security}

Security is intentionally limited in the current design because, as discussed in \cref{ds-features}, it is not a primary driver of this project. The system is assumed to run in controlled laboratory, development, or trusted demonstration environments.

Accordingly, the current design does \textbf{not} make authentication a core architectural dependency for either peer-to-peer traffic or read-only observation endpoints. Similarly, it does not define a rich authorization model with differentiated user permissions. The effective roles are simple: nodes participate in the distributed protocol, while observers query read-only information.

This choice keeps the architecture aligned with the educational scope of the project and avoids obscuring the main distributed-systems concerns with security infrastructure that is orthogonal to the report goals. However, the design has been kept modular enough that stronger security policies could be introduced later at clear boundaries:
\begin{itemize}
  \item on node-to-node channels, to authenticate peers and protect inter-node communication;
  \item on HTTP observation endpoints, to restrict access to monitoring data;
  \item on deployment boundaries, through network isolation and reverse-proxy policies.
\end{itemize}

Likewise, no cryptographic protocol is a conceptual prerequisite of the current system model. Confidentiality and strong identity guarantees are therefore outside the present scope and remain future extensions rather than core design commitments.

\clearpage
